\chapter{Đánh giá hệ thống}
\textit{Chương này trình bày quá trình đánh giá hệ thống dựa trên các tiêu chí về chức năng và phi chức năng, qua đó làm rõ mức độ đáp ứng của hệ thống đối với các mục tiêu đã đề ra.}

\section{Mục tiêu đánh giá}
Quá trình đánh giá hệ thống nhắm đến hai nhóm mục tiêu chính:

\begin{itemize}
    \item \textbf{Yêu cầu chức năng (Functional Requirements):} kiểm chứng rằng tất cả các chức năng của hệ thống hoạt động đúng theo đặc tả, bao gồm các luồng nghiệp vụ cốt lõi như đăng nhập, đăng ký, làm bài thi, quản lý nội dung và các tính năng AI.
    \item \textbf{Yêu cầu phi chức năng (Non-functional Requirements):} đo lường hiệu năng, khả năng chịu tải và tính nhất quán của hệ thống trên các môi trường triển khai khác nhau, nhằm đảm bảo hệ thống đáp ứng các tiêu chí chất lượng đã cam kết.
\end{itemize}

\section{Tiêu chí và phương pháp đánh giá}

\subsection{Kiểm thử yêu cầu chức năng --- Kiểm thử End-to-End}

Toàn bộ yêu cầu chức năng của hệ thống được kiểm chứng thông qua bộ kiểm thử End-to-End (E2E) sử dụng \textbf{Playwright} --- một framework kiểm thử tự động hóa trình duyệt. Playwright cho phép mô phỏng các thao tác thực tế của người dùng trên giao diện web, bao gồm điều hướng, nhập liệu, nhấn nút và kiểm tra kết quả hiển thị, qua đó kiểm chứng toàn bộ luồng nghiệp vụ từ đầu đến cuối.

Bộ test case E2E được tổ chức thành hai nhóm tương ứng với hai hệ thống: \textbf{English Prep} và \textbf{Lingriser}, bao phủ các nhóm chức năng chính gồm xác thực, quản lý hồ sơ, quản lý bài thi, flashcard, chat, blog, thông báo và quản trị hệ thống. Toàn bộ danh sách test case được trình bày chi tiết trong phần Phụ lục.

\subsection{Kiểm thử yêu cầu phi chức năng}

Các yêu cầu phi chức năng được kiểm thử với phương pháp và công cụ tương ứng như sau:

\begin{table}[H]
\centering
\caption{Phương pháp kiểm thử yêu cầu phi chức năng}
\renewcommand{\arraystretch}{1.4}
\small
\begin{tabular}{|p{0.08\linewidth}|p{0.22\linewidth}|p{0.25\linewidth}|p{0.35\linewidth}|}
\hline
\textbf{ID} & \textbf{Tên NFR} & \textbf{Công cụ} & \textbf{Phương pháp} \\
\hline
NFR1 & Performance & Postman, k6 & Đo thời gian phản hồi API bằng Postman cho các endpoint đơn lẻ. Sử dụng k6 để chạy kịch bản tải với nhiều mức virtual user (VU), đo response time trung bình và tỷ lệ request thành công. \\
\hline
NFR2 & Scalability & k6 & Chạy stress test với k6, tăng dần số lượng concurrent user lên đến 1.000 VU, đo khả năng xử lý đồng thời và tính ổn định của hệ thống dưới tải cao. \\
\hline
NFR4 & Portability & CI/CD Pipeline & Kiểm chứng thông qua pipeline CI/CD: cùng một bộ mã nguồn được build, test và chạy thành công trên cả môi trường Development và Production, đảm bảo hành vi nhất quán giữa các môi trường. \\
\hline
\end{tabular}
\end{table}

\section{Kết quả đánh giá}

\subsection{Kết quả kiểm thử E2E}

Toàn bộ bộ kiểm thử E2E được thực thi trên cả hai hệ thống English Prep và Lingriser. Kết quả cho thấy \textbf{100\% test case đều pass}, xác nhận rằng tất cả các chức năng chính của hệ thống hoạt động đúng theo đặc tả trong điều kiện vận hành bình thường.

Tuy nhiên, khi hệ thống chịu tải cao (stress test), một số request có thể bị từ chối do cơ chế \textbf{throttle guard} --- bộ giới hạn tốc độ request được tích hợp nhằm bảo vệ hệ thống khỏi bị quá tải. Đây là hành vi có chủ đích trong thiết kế hệ thống: throttle guard ưu tiên duy trì tính ổn định tổng thể thay vì cho phép tất cả request đi qua khi tài nguyên không đủ đáp ứng.

\subsection{Kết quả kiểm thử hiệu năng với k6}

Bảng~\ref{tab:k6-results} tổng hợp kết quả chạy stress test bằng k6 trên ba nhóm kịch bản khác nhau, từ các endpoint công khai không yêu cầu xác thực đến kiểm thử toàn tải với 1.000 virtual user đồng thời.

\begin{table}[H]
\centering
\caption{Kết quả kiểm thử hiệu năng với k6}
\label{tab:k6-results}
\renewcommand{\arraystretch}{1.4}
\small
\begin{tabular}{|p{0.2\linewidth}|c|c|r|c|c|c|}
\hline
\textbf{Kịch bản} & \textbf{Thời gian} & \textbf{Max VU} & \textbf{Tổng request} & \textbf{Check pass} & \textbf{Avg response} & \textbf{JSON output} \\
\hline
Public endpoints & 4 phút & 10 & 944 & 100\% & 5,71 ms & 3,7 MB \\
\hline
Auth-required endpoints & 18 phút & 1.000 & 237.253 & 98,35\% & 1,8 ms & 988 MB \\
\hline
Full Stress & 18 phút & 1.000 & 246.301 & 98,83\% & 5,86 ms & 1,1 GB \\
\hline
\end{tabular}
\end{table}

\noindent Phân tích kết quả:

\begin{itemize}
    \item \textbf{Public endpoints:} với 10 virtual user, hệ thống xử lý 944 request trong 4 phút với tỷ lệ thành công 100\% và thời gian phản hồi trung bình 5,71ms --- cho thấy các endpoint công khai hoạt động ổn định và đáp ứng nhanh.
    \item \textbf{Auth-required endpoints:} khi tăng lên 1.000 VU đồng thời trong 18 phút, hệ thống xử lý hơn 237.000 request với tỷ lệ check pass đạt 98,35\%. Thời gian phản hồi trung bình chỉ 1,8ms cho thấy cơ chế xác thực và xử lý nghiệp vụ hoạt động hiệu quả. Tỷ lệ 1,65\% request không pass chủ yếu do throttle guard kích hoạt khi lượng request vượt ngưỡng cho phép.
    \item \textbf{Full Stress:} kịch bản tải tổng hợp với 1.000 VU và hơn 246.000 request đạt tỷ lệ check pass 98,83\%. Thời gian phản hồi trung bình 5,86ms vẫn nằm trong ngưỡng chấp nhận được, xác nhận hệ thống có khả năng chịu tải ở mức cao mà không suy giảm hiệu năng đáng kể.
\end{itemize}

\noindent Đối chiếu với các yêu cầu phi chức năng:

\begin{itemize}
    \item \textbf{NFR1 --- Performance:} thời gian phản hồi trung bình dao động từ 1,8ms đến 5,86ms, thấp hơn đáng kể so với ngưỡng 500ms đã đặt ra cho các thao tác non-AI. Yêu cầu này được đáp ứng.
    \item \textbf{NFR2 --- Scalability:} hệ thống xử lý thành công hơn 246.000 request với 1.000 VU đồng thời, tỷ lệ thành công trên 98\%. Mặc dù chưa đạt đến mốc 10.000 concurrent user do giới hạn tài nguyên VPS hiện tại, kết quả cho thấy kiến trúc microservices cho phép mở rộng khi bổ sung thêm tài nguyên.
    \item \textbf{NFR4 --- Portability:} pipeline CI/CD chạy thành công trên cả hai môi trường Development và Production với cùng bộ mã nguồn, xác nhận tính nhất quán giữa các môi trường triển khai.
\end{itemize}

\section{Kết luận chương}
Chương này đã trình bày quá trình đánh giá hệ thống dựa trên cả yêu cầu chức năng và phi chức năng. Kết quả kiểm thử E2E cho thấy toàn bộ chức năng hoạt động đúng đặc tả với tỷ lệ pass 100\%. Về hiệu năng, hệ thống đáp ứng tốt các chỉ tiêu đã đề ra với thời gian phản hồi trung bình dưới 6ms và tỷ lệ thành công trên 98\% ngay cả khi chịu tải cao 1.000 virtual user. Các kết quả này khẳng định hệ thống đã sẵn sàng cho môi trường vận hành thực tế, đồng thời chỉ ra hướng cải thiện về khả năng mở rộng khi nhu cầu sử dụng tăng cao.
